\documentclass[10pt,a4paper]{article}
\usepackage{amsmath,amssymb,amsthm,graphicx}
\newcommand{\bits}{\{0,1\}}
\newcommand{\DT}{\operatorname{DT}}
\newcommand{\DTSize}{\operatorname{DTSize}}
\newcommand{\CC}{\operatorname C}
\newcommand{\comp}{\diamond}
\newcommand{\OR}{\operatorname{OR}}
\newcommand{\AND}{\operatorname{AND}}
\newcommand{\XOR}{\operatorname{XOR}}
\newenvironment{originalresult}[1]{\par\medskip\noindent\textbf{#1.}\itshape}{\par\medskip}
\title{Complete gadget dichotomy for deterministic communication lifting}
\author{Yaroslav Alekseev, Yuval Filmus, Alexander V. Smal}
\date{doi:10.1007/s00037-025-00276-5; CC BY 4.0}
\begin{document}
\maketitle
\noindent\textit{Editorial note.} Faithful editable transcription of actual complete CC BY4.0 journal version. Full author Theorem6.9 proof and every source-local certificate/matrix lemma it uses are retained. All source equations and restrictions are preserved; no mathematical proof is supplemented or repaired. Original AND/OR/XOR matrices are editable displayed formulas. Standard cited communication and parity-query facts remain precise prerequisites.\par\medskip
\noindent\textit{Allowed standard prerequisites (editorial).}\begin{itemize}
\item Randomized unique set-disjointness lower bound. Razborov1992; exact use retained in full author reduction
\item Deterministic equality lower bound. Kushilevitz/Nisan1996, exact invoked statement retained
\item XOR communication to parity-query polynomial bound. Hatami/Hosseini/Lovett2018; exact bound and source Corollary6.11 retained
\item Parity certificate/decision-tree polynomial relations. Zhang/Shi2010; exact source relations retained
\end{itemize}\medskip
\section*{Block composition (author source p.2)}
For $f:\bits^n\to V$ and $g:\bits^m\to\bits$, a (block-)composition $f\comp g:\bits^{n\times m}\to V$ is defined by
\[(f\comp g)(z_1,z_2,\ldots,z_n):=f(g(z_1),g(z_2),\ldots,g(z_n)),\]
where each $z_i\in\bits^m$.
\section*{Prerequisites (author source pp.10--12)}
All the logarithms are base 2. For $g:\bits^k\to\bits$, we define a function $g^n:\bits^{n\times k}\to\bits^n$ such that
\[g^n(x_1,x_2,\ldots,x_n):=(g(x_1),g(x_2),\ldots,g(x_n)).\]
For a function $f:\bits^n\to\bits$, an input $x\in\bits^n$, and a partial assignment $\sigma\in\{0,1,*\}^n$ that agrees with $x$ on all non-$*$ coordinates, we use $f|_\sigma$ to denote a restriction of $f$ to $\sigma$ and $x|_\sigma$ to denote a projection of $x$ to the $*$-coordinates of $\sigma$.
\paragraph{Decision tree complexity.} A decision tree is a rooted binary tree with internal nodes labeled by input variables (representing queries), edges labeled with 0 and 1, and leaves labeled by some values. Decision tree evaluation is defined in a natural way: given an assignment for the input variables, we traverse the tree starting from the root and then sequentially choose every next edge according to the given assignments until we reach some leaf. The label of this leaf is the result of the evaluation. A decision tree computes some function $f$ if for all possible assignments to the input variables, the decision tree evaluates to the value of the function. For a function $f$, $\DT(f)$ denotes the minimal depth of a decision tree computing $f$, and $\DTSize(f)$ denotes the minimal number of leaves in a decision tree computing $f$.
\paragraph{Certificate complexity.} A certificate complexity is a non-deterministic analogue of the decision tree complexity. A certificate for an input $x\in\bits^n$ to a function $f$ is a set $S\subseteq[n]$ of indices such that $f$ restricted to all inputs that match $x$ on $S$ is constant, i.e., $f(y)=f(x)$ whenever $y|_S=x|_S$. The certificate complexity $\CC(f,x)$ of $f$ at input $x$ is the size of the smallest certificate for $x$. Finally, the certificate complexity of function $f$ is defined as $\CC(f):=\max_{x\in\bits^n}\CC(f,x)$.

\section*{Complexity definitions (author source pp.11--13)}
\paragraph{Parity decision tree complexity.} A parity decision tree extends the concept of a decision tree by allowing queries that are arbitrary linear combinations of the input variables. It can be described as a rooted binary tree with internal nodes labeled by linear combinations of the input variables, edges labeled with 0 and 1, and leaves labeled by some values. The evaluation is defined analogously. For a function $f$, $\DT_\oplus(f)$ denotes the minimal depth of a parity decision tree computing $f$, and $\DTSize_\oplus(f)$ denotes the minimal number of leaves in a parity decision tree computing $f$.

\paragraph{Deterministic communication complexity and rank.} For a function $f:X\times Y\to Z$, consider the following game for two players, Alice and Bob, that want to compute $f$. Alice and Bob are given $x\in X$ and $y\in Y$, respectively. Their goal is to compute $f(x,y)$. In order to do it, the players exchange information about their parts of the input using a simple communication channel that allows sending bit messages. Before the game, the players come up with a communication protocol that determines their behavior on all possible inputs. The cost of a protocol is the maximum total number of bits sent by the players over all possible inputs $x\in X,y\in Y$. The deterministic communication complexity of $f$ is the minimal cost of a deterministic communication protocol that computes $f$. We denote it by $D(f)$.
For a function $f:X\times Y\to\bits$, a communication matrix of $f$ is a matrix $M_f\in\bits^{X\times Y}$ such that $(M_f)_{x,y}:=f(x,y)$ for all $x\in X,y\in Y$. We define $\operatorname{rank}(f)$ to be the rank of the matrix $M_f$ over the reals.

\section*{Block sensitivity and randomized communication (author source pp.12--13)}
A block $B\subseteq[n]$ is sensitive for a function $f$ at $x$ iff $f(x)\ne f(x^B)$. The block sensitivity $\operatorname{bs}(f,x)$ of $f$ at $x$ is the maximum number of disjoint sensitive blocks for $f$ at $x$. The block sensitivity $\operatorname{bs}(f)$ of $f$ is the maximum sensitivity $\operatorname{bs}(f,x)$ over all points $x\in\bits^n$.
We use the notation $x^B$ for $x\in\bits^n$ and $B\subseteq[n]$ to denote $x$ with all the coordinates in $B$ flipped.
In a randomized communication game, Alice and Bob have access to an unlimited amount of random bits, and they can use it to decide which bit to send next. We say that a (private coin) randomized communication protocol $\Pi$ computes a function $f:X\times Y\to Z$ with error $\varepsilon$ if
\[\Pr_{r_A,r_B}[\Pi(x,y,r_A,r_B)=f(x,y)]\ge1-\varepsilon,\qquad\forall x\in X,y\in Y,\]
where $r_A$ and $r_B$ are the strings of random bits used by Alice and Bob, respectively. The cost of a randomized protocol is the maximum number of bits that can be sent. We denote by $R_\varepsilon(f)$ the minimal cost of a private coin randomized communication protocol that computes $f$ with error $\varepsilon$. The randomized communication complexity of $f$ is defined as $R(f):=R_{1/3}(f)$.
\section*{6.1 Reductions in communication complexity (author source p.35)}
Consider a pair of two-party functions $f_i:X_i\times Y_i\to\bits$ for $i=1,2$. We say that $f_1$ reduces to $f_2$, denoted $f_1\le f_2$, if the communication matrix of $f_1$ is a submatrix of the communication matrix of $f_2$. Equivalently, $f_1\le f_2$ iff there exist one-to-one mappings $\pi_A$ and $\pi_B$ such that
\[f_1(x,y)=f_2(\pi_A(x),\pi_B(y)),\qquad\forall(x,y)\in X_1\times Y_1.\]

\section*{Parity certificates (author source pp.12,24)}
A parity certificate for an input $x\in\bits^n$ to a function $f$ is a set $A$ of linear forms such that $f$ is constant on the affine subspace defined by $A=A(x)$. A parity certificate complexity $C_\oplus(f,x)$ of a function $f$ at input $x$ is the size (number of linear forms) of the smallest parity certificate for $x$. The parity certificate complexity of $f$ is defined as $C_\oplus(f):=\max_{x\in\bits^n}C_\oplus(f,x)$.
\section*{5.2 Minimum weight lemma}
A parity certificate $C$ is a system of affine equations over input variables. It can be represented by a matrix $A$ and a vector $b$ that define a linear subspace, i.e., $Ax=b$.
\begin{originalresult}{Definition 5.2}
The minimum weight of a parity certificate $C$ is the minimum Hamming weight of nonzero vectors in the row space of $A$.
\end{originalresult}
For example, one can consider a certificate $\{x+y=1,x+y+z=0\}$. The minimum weight of this certificate is equal to 1, and this corresponds to the equation $z=1$. On the other hand, the minimum weight of a certificate $\{x+y=1,x+z=0\}$ is equal to 2.

\begin{originalresult}{Lemma 5.3 (minimal weight lemma; author source p.25)}
For $f:\bits^n\to\bits$ and $g:\bits^m\to\bits$, suppose that $f\comp g$ has parity certificate complexity at most $k$ at some point $x$. Let $K$ be a parameter. If the certificate complexity of $f$ at $g^n(x)$ is greater than $k\cdot K$ then we can find a blockwise partial assignment $\sigma$ to the inputs of $f\comp g$ consistent with $x$, such that for some $k'\le k$:
\begin{itemize}
\item $(f\comp g)|_\sigma$ has a parity certificate of size $k'$ at $x|_\sigma$ whose minimum weight is at least $K$,
\item $C(f|_{g^n(\sigma)},g^n(x)|_{g^n(\sigma)})>k'K$. (Note that $\sigma$ is blockwise and agrees with $x$, so here we consider the certificate complexity of $f$ at $g^n(x)$ on all the blocks that are not set by $\sigma$).
\end{itemize}
\end{originalresult}
\begin{proof}
The proof is by induction on $k$. If $k=0$, then the premises contradict each other, so we can assume that $k>0$.
Let $C$ be the parity certificate for $f\comp g$ at input $x$. By the assumption, the size of $C$ is at most $k$. If its minimum weight is at least $K$, then we are done. Otherwise, without loss of generality, $C$ contains a nontrivial constraint $\ell=\alpha$ whose support is less than $K$. Choose a partial assignment that agrees with $x$ on the support of this constraint, and complete it to a blockwise partial assignment $\sigma$ consistent with $x$ (complete only the blocks with at least one variable assigned). Also, let $\tau:=g^n(\sigma)$.
By construction, $(f\comp g)|_\sigma$ has a parity certificate of size at most $k-1$ at $x|_\sigma$, namely $C\setminus\{\ell=\alpha\}$. On the other hand, since we only set the values of at most $K$ blocks, the certificate complexity of $f$ reduces by at most $K$ after being restricted to $\tau$, and hence $C(f|_\tau,g^n(x)|_\tau)>(k-1)K$. Now we can apply the induction hypothesis for $k-1$.
\end{proof}
The following lemma demonstrates an essential property of the parity certificates with huge minimal weight:
\begin{originalresult}{Lemma 5.4 (author source pp.25--26)}
Suppose that $C_\oplus$ is a parity certificate with minimum weight at least $K$ and $C$ is a (regular) certificate of size less than $K$. Then there is a point that satisfies both $C$ and $C_\oplus$.
\end{originalresult}
In the proof, the parity certificate is processed with a careful Gaussian elimination, and then the variables fixed by the certificate $C$ are free variables of the parity certificate.
\begin{proof}
Let $V+b$ be the affine subspace corresponding to $C_\oplus$, where $V$ is a linear subspace. We pick a basis $\ell_1,\ldots,\ell_k$ for $V$ and indices $i_1,\ldots,i_k\in[n]$ such that:
\begin{itemize}
\item $z_{i_t}$ is in the support of $\ell_t$,
\item $z_{i_s}$ is not in the support of $\ell_t$ for $s<t$,
\item $z_{i_t}$ does not belong to the support of $C$.
\end{itemize}
We do it inductively. At step $t$, we have already picked $\ell_1,\ldots,\ell_{t-1}$. Consider vectors in $V$ whose support does not contain any of the variables $z_{i_1},\ldots,z_{i_{t-1}}$. Since $V$ has minimum weight $K$, we can find such a vector $\ell_t$ whose support contains a variable $z_{i_t}$ with $z_{i_t}\notin C$. By the assumption, this vector does not contain any of the variables $z_{i_1},\ldots,z_{i_{t-1}}$.
Now we can construct an input $\alpha$ which satisfies both $C$ and $C_\oplus$. First, we set the variables from $C$ such that the resulting input satisfies $C$. Then we set arbitrarily all the variables not occurring in $C$ or $\{z_{i_1},\ldots,z_{i_k}\}$. Finally, we set variables $z_{i_k},\ldots,z_{i_1}$ (in this order) in such a way that the resulting input satisfies $C_\oplus$.
\end{proof}

\begin{originalresult}{Theorem 6.1 (Zhang; author source p.36)}
If a two-party gadget $g:X\times Y\to\bits$ satisfies $\AND,\OR\le g$, then for every function $f:\bits^n\to Q$, the function $f\comp g$ has (constant error) randomized communication complexity $\Omega(\operatorname{bs}(f))$.
\end{originalresult}
For completeness, we present the proof that is a reduction from unique set disjointness.
\begin{originalresult}{Definition 6.2}
The disjointness function $\operatorname{Disj}_n:\bits^n\times\bits^n$ is defined as follows:
\[\operatorname{Disj}_n(x,y):=\bigvee_{i=1}^n x_i\wedge y_i.\]
Unique set disjointness, $\operatorname{PromiseDisj}_n$, is a promise version of disjointness with the promise that there is always at most one $i$ such that $x_i\wedge y_i=1$.
\end{originalresult}
\begin{proof}[Proof of Theorem 6.1]
Razborov (1992) showed that $R(\operatorname{PromiseDisj}_m)=\Omega(m)$. We will show a reduction from this problem to computing $f\comp g$. To illustrate the proof, let us start with the case in which $f(0)=0$ and $f(e_i)=1$ for all $i\in[n]$, where $e_i$ is the $i$'th unit vector. Then
\[\operatorname{PromiseDisj}_n(x,y)=f(x_1\wedge y_1,\ldots,x_n\wedge y_n).\]
Since $\AND\le g$, we can solve unique set disjointness using a protocol for $f\comp g$. Therefore, $f\comp g$ has communication complexity $\Omega(n)$. Note that this case only required $\AND\le g$; we need $\OR\le g$ in order to handle inputs changing from 1 to 0.
To handle the general case, let inputs $X^\wedge(0),X^\wedge(1)\in X$ and $Y^\wedge(0),Y^\wedge(1)\in Y$ witness $\AND\le g$, and let $X^\vee(0),X^\vee(1)\in X$ and $Y^\vee(0),Y^\vee(1)\in Y$ witness $\OR\le g$.
Let $z$ be a point witnessing $\operatorname{bs}(f)$, and let $B_1,\ldots,B_m$ be the corresponding blocks, where $m=\operatorname{bs}(f)$. Given inputs $x,y$ to unique set disjointness, we construct the following inputs $X,Y$ to $f\comp g$:
\begin{itemize}
\item If $i$ does not belong to any of the blocks, then we fix $X_i,Y_i$ so that $g(X_i,Y_i)=z_i$.
\item If $i\in B_j$ and $z_i=0$, put $X_i=X^\wedge(x_j),Y_i=Y^\wedge(y_j)$.
\item If $i\in B_j$ and $z_i=1$, put $X_i=X^\vee(\bar x_j),Y_i=Y^\vee(\bar y_j)$.
\end{itemize}
The reader can check that if $x,y$ are disjoint, then applying $g$ to $X,Y$ yields $z$, and if they intersect only at $j$, then it yields $z\oplus B_j$. Hence, $f\comp g$ can be used to compute unique set disjointness of length $m$, and consequently, we obtain a lower bound of $\Omega(m)$.
\end{proof}

\section*{Block-diagonal gadgets (author source pp.37--38)}
For the classification of gadgets, we will need the following lemma that shows that if $\OR\nleq g$ then the communication matrix of $g$ is block diagonal.
\begin{originalresult}{Lemma 6.3}
If $g:X\times Y\to\bits$ satisfies that $\OR\nleq g$ then there are partitions of $X$ and $Y$ into disjoint sets
\[X=X_0\sqcup X_1\sqcup\cdots\sqcup X_k\quad\text{and }Y=Y_0\sqcup Y_1\sqcup\cdots\sqcup Y_k,\]
such that
\begin{itemize}
\item If $x\in X_i,y\in Y_j$, where $i\ne j$, or $i=0$, or $j=0$, then $g(x,y)=0$.
\item If $x\in X_i,y\in Y_i$, where $i\ne0$, then $g(x,y)=1$.
\end{itemize}
\end{originalresult}
\begin{proof}
We can construct such a partition step by step. Consider any point $(x_0,y_0)$ such that $g(x_0,y_0)=1$. If we don't have such a point, then we can take $X_0=X$ and $Y_0=Y$. Otherwise, we consider all points $y\in Y$ such that $g(x_0,y)=1$. Denote the set of all such points as $Y'$. Now, we can consider all $x\in X$ such that $g(x,y_0)=1$. Denote the set of all such points as $X'$. Observe that every point $(x,y)\in X'\times Y'$ satisfies $g(x,y)=1$, otherwise we can project our gadget to points $(x_0,y_0),(x_0,y),(x,y_0),(x,y)$ and get OR.
Now, observe that for any $(x'',y')\in(X\setminus X')\times Y'$ we have $g(x'',y_0)=0$ by construction, and thus $g(x'',y')=0$, since otherwise we can project the gadget to points $(x_0,y'),(x_0,y_0),(x'',y'),(x'',y_0)$ and get OR. By analogy, for any $(x',y'')\in X'\times(Y\setminus Y')$ we have $g(x',y'')=0$.
This means that $X'\times Y'$ forms a valid block of our partition, so we can use the same procedure to construct a block of $(X\setminus X')\times(Y\setminus Y')$. In the end, if we get sets $X''$ and $Y''$ such that any $(x,y)\in X''\times Y''$ satisfies $g(x,y)=0$, we take $X_0=X''$ and $Y_0=Y''$.
\end{proof}

\section*{Blow-up classification (author source pp.38--39)}
\begin{originalresult}{Definition 6.4}
We say that a gadget $g:X\times Y\to\bits$ is a blow-up of a gadget $h:Z\times W\to\bits$ if there are decompositions $X=\bigsqcup_{z\in Z}X_z$ and $Y=\bigsqcup_{w\in W}Y_w$, with $X_z,Y_w\ne\varnothing$, such that all $(x_z,y_w)\in X_z\times Y_w$ satisfy $g(x_z,y_w)=h(z,w)$.
\end{originalresult}
For example, $g$ is a blow-up of XOR if (up to rearrangement) it has a communication matrix of the following form
\[\begin{bmatrix}1&\cdots&1&0&\cdots&0\\\vdots&\ddots&\vdots&\vdots&\ddots&\vdots\\1&\cdots&1&0&\cdots&0\\0&\cdots&0&1&\cdots&1\\\vdots&\ddots&\vdots&\vdots&\ddots&\vdots\\0&\cdots&0&1&\cdots&1\end{bmatrix}.\]
If $g$ is a blow-up of $h$, then $f\comp g$ and $f\comp h$ have the same communication complexity in all models. Indeed, on the one hand, $h$ is a restriction of $g$, and so a protocol for $f\comp g$ can be used to solve $f\comp h$; and on the other hand, by replacing $x\in X_z$ by $z$ and $y\in Y_w$ by $w$, we can use a protocol for $f\comp h$ to solve $f\comp g$.
\begin{originalresult}{Corollary 6.5}
If $g:X\times Y\to\bits$ satisfies both $\OR\nleq g$ and $\AND\nleq g$ then either $g$ is constant, or it is a blow-up of XOR.
\end{originalresult}
\begin{proof}
Let $X=X_0\sqcup X_1\sqcup\cdots\sqcup X_k$ and $Y=Y_0\sqcup Y_1\sqcup\cdots\sqcup Y_k$ be the decomposition provided by Lemma 6.3. We can assume that the $X_i,Y_i\ne\varnothing$, except that possibly $X_0=Y_0=\varnothing$.
If $X_0\ne\varnothing$ and $k>0$ then $\AND\le g$. Indeed, if we restrict the matrix of $g$ so that it contains one row each from $X_0,X_1$ and one column each from $Y_0,Y_1$ then we get
\[\begin{bmatrix}0&0\\0&1\end{bmatrix}\]
which is AND. Consequently, either $k=0$ or $X_0=Y_0=\varnothing$.
If $k=0$ then $g\equiv0$. Now suppose that $X_0=Y_0=\varnothing$. If $k\ge3$ and we restrict the communication matrix of $g$ so that it contains one row from each $X_1,X_2,X_3$ and one column from each $Y_1,Y_2,Y_3$ then we get
\[\begin{bmatrix}\mathbf1&0&\mathbf0\\\mathbf0&1&\mathbf0\\0&0&1\end{bmatrix}.\]
Again, we get a restriction to AND (the entries in bold text correspond to the AND matrix). Hence, $k\le2$. If $k=1$ then $g\equiv1$, and otherwise $g$ is a blow-up of XOR.
\end{proof}
\begin{originalresult}{Corollary 6.6}
If $g:X\times Y\to\bits$ satisfies both $\OR\nleq g$ and $\XOR\nleq g$ then either $g$ is constant, or it is a blow-up of AND. Similarly, if $\AND\nleq g$ and $\XOR\nleq g$ then either $g$ is constant, or it is a blow-up of OR.
\end{originalresult}
\begin{proof}
We can assume that the $X_i,Y_i\ne\varnothing$, except that possibly $X_0=Y_0=\varnothing$. We start by proving the first claim, and so we assume that $\OR,\XOR\nleq g$.
If $k\ge2$ then $\XOR\le g$. Indeed, if we restrict the matrix of $g$ so that it contains one row from each $X_1,X_2$ and one column from each $Y_1,Y_2$ then we obtain
\[\begin{bmatrix}1&0\\0&1\end{bmatrix}\]
which is XOR. If $k=0$ then $g\equiv0$. If $k=1$ and $X_0=Y_0=\varnothing$ then $g\equiv1$. Finally, if $k=1$ and $X_0,Y_0\ne\varnothing$ then $g$ is a blow-up of AND.
Suppose that $\AND,\XOR\nleq g$. Then $\OR,\XOR\nleq\bar g$, where $\bar g$ is the negation of $g$. Hence, either $\bar g$ is constant or it is a blow-up of AND. Thus, either $g$ is constant or it is a blow-up of OR.
\end{proof}

\section*{6.3 Deterministic communication lifting support (author source pp.41--45)}
\begin{originalresult}{Theorem 6.9}
If a two-party gadget $g:X\times Y\to\bits$ satisfies $\AND,\XOR\le g$, then for every function $f:\bits^n\to\bits$, the function $f\comp g$ has communication complexity $\Omega(\operatorname{bs}(f)^k)$ for some fixed constant $k>0$. The same is true if $g$ satisfies $\OR,\XOR\le g$.
\end{originalresult}
We start with the following observation: according to Hatami et al.(2018), for any $f:\bits^n\to\bits$ we have $\DT_\oplus(f)\le O(D(f\comp\XOR)^6)$. Zhang \& Shi(2010) showed that $\DT_\oplus(f)=O(C_\oplus(f)^2)$. Combining these results, we can derive the following corollary:
\begin{originalresult}{Corollary 6.11}
For any $f:\bits^n\to\bits$, we have
\[C_\oplus(f)\le O(D(f\comp\XOR)^{12}).\]
\end{originalresult}
Using this corollary, our main goal will be to find a substitution $\tau$ and two points $x$ and $y$ such that $f(x)\ne f(y)$ and $x|_\tau$ and $y|_\tau$ have parity certificates of large minimal weight. Then, using Lemma6.12 and Lemma6.13, we will get the required lifting.
In the following lemma, a minimal sensitive block of a function $f:\bits^n\to\bits$ at point $x$ is a subset $I\subseteq[n]$ such that $f(x^I)\ne f(x)$ but $f(x^J)=f(x)$ for all $J\subsetneq I$. This lemma shows that the deterministic communication complexity of $f\comp\XOR$ is lower bounded by the size of any minimal sensitive block of $f$.
\begin{originalresult}{Lemma 6.12}
For any $f:\bits^n\to\bits$ and $\alpha\in\bits^n$ such that there is a minimal sensitive block of $f$ at $\alpha$ of size $S$, the following holds:
\[D(f\comp\XOR)\ge S.\]
\end{originalresult}
\begin{proof}
Consider some minimal sensitive block of $f$ at $\alpha$ of size $S$. Let $i_1,i_2,\ldots,i_S$ be the indices of the variables in this block of $\alpha$. Let us restrict the function $f$ by substituting values $x_j=\alpha_j$ for all $j\notin\{i_1,i_2,\ldots,i_S\}$. This gives us a function $f'(y_1,y_2,\ldots,y_S)$ such that
\begin{itemize}
\item $f'(y_1,\ldots,y_S)=f(\alpha)$, if $y_j=\alpha_{i_j}$ for at least one $j$.
\item $f'(y_1,\ldots,y_S)\ne f(\alpha)$, if $y_j\ne\alpha_{i_j}$ for all $j$.
\end{itemize}
Thus, if we substitute XOR inside each $y_j$ this will give us an instance of the function
\[(z_1\oplus z_1'\oplus\alpha_{i_1})\vee(z_2\oplus z_2'\oplus\alpha_{i_2})\vee\cdots\vee(z_S\oplus z_S'\oplus\alpha_{i_S}),\]
where each $z_i$ and $z_i'$ are distinct variables. This function is essentially a negation of the equality function with a shift. So, its deterministic complexity equals $S$ (see Kushilevitz \& Nisan(1996) for the lower bound for $\operatorname{EQ}_n$). Thus,
\[D(f\comp\XOR)\ge D(f'\comp\XOR)\ge S.\]
\end{proof}
The following lemma shows that the deterministic communication complexity of $f\comp\AND$ is lower bounded by the block sensitivity of $f$ at the all-0 input.
\begin{originalresult}{Lemma 6.13}
For any $f:\bits^n\to\bits$ such that $\operatorname{bs}(f,0)=S$, the following holds:
\[D(f\comp\AND)\ge\Omega(S).\]
\end{originalresult}
\begin{proof}
The proof of the lemma essentially follows the proof from Zhang(2009). Let $I_1,\ldots,I_S$ be the minimal sensitive blocks of $f$ at 0. We describe a reduction from $\operatorname{PromiseDisj}_S$ to $f\comp\AND$, so any protocol for $f\comp\AND$ can be used to solve $\operatorname{PromiseDisj}_S$, and thus $D(f\comp\AND)\ge\Omega(S)$. For any input $(a,b)\in\bits^S$ to $\operatorname{PromiseDisj}_S$ we define the corresponding input $(x,y)\in\bits^n$ to $f\comp\AND$.
\begin{enumerate}
\item Let $x_i=y_i=0$ for all $i\notin\bigcup_{j=1}^S I_j$.
\item Let $x_i=a_j$ and $y_i=b_j$ for all $i\in I_j$.
\end{enumerate}
We claim that this is a correct reduction. Indeed, if there is a unique $j$ such that $a_j\wedge b_j=1$, then
\[f(x_1\wedge y_1,x_2\wedge y_2,\ldots,x_n\wedge y_n)=f(0^{I_j})\ne f(0,\ldots,0).\]
On the other hand, if $a_j\wedge b_j=0$ for all $j\in[S]$, then
\[f(x_1\wedge y_1,x_2\wedge y_2,\ldots,x_n\wedge y_n)=f(0,\ldots,0).\]
\end{proof}
The next lemma is an application of the `minimum weight lemma' for the identity gadget.
\begin{originalresult}{Lemma 6.14}
Suppose that $f$ has a parity certificate at some point $x$ of size at most $k$. Let $K$ be a parameter. Suppose that the certificate complexity of $f$ at $x$ is greater than $k\cdot K$. Then we can find a partial assignment $\sigma$ to at most $k\cdot K$ variables of $f$, consistent with $x$, such that:
\begin{itemize}
\item $f|_\sigma$ has a parity certificate at $x|_\sigma$ of size at most $k$.
\item The minimum weight of this certificate is at least $K$.
\end{itemize}
\end{originalresult}
\begin{proof}
Apply Lemma5.3 for the function $f$ and the identity function as a gadget $g$.
\end{proof}
For the last ingredient needed for the proof of the classification theorem, the following lemma shows that parity certificates of high minimum weight always intersect with (regular) certificates of sufficiently small size.
\begin{originalresult}{Lemma 6.15}
Suppose that for some function $f$ and some input $x$ the following holds: $f$ has a parity certificate $C_\oplus(x)$ of size $k$ at $x$ whose minimum weight is at least $K$. Then, for any point $y$ with certificate complexity less than $K$, there is a non-empty intersection of the corresponding certificates $C(y)$ and $C_\oplus(x)$.
\end{originalresult}
\begin{proof}
Apply Lemma5.4 with $C_\oplus=C_\oplus(x)$ and $C=C(y)$.
\end{proof}
Now we have everything we need to prove Theorem6.9.
\begin{proof}[Theorem 6.9]
Assume that $\AND,\XOR\le g$. The proof in the case of $\OR,\XOR\le g$ is similar. Consider some function $f$. There is a value $S$ such that $\operatorname{bs}(f)=9S^5$. If $C_\oplus(f)\ge S$, then by Corollary6.11,
\[D(f\comp g)\ge D(f\comp\XOR)\ge C_\oplus(f)^{1/12}\ge S^{1/12}\ge\Omega(\operatorname{bs}(f)^{1/(12\cdot5)}),\]
and so the theorem holds for $k=1/60$. So, we can assume that $C_\oplus(f)<S$. Also, we can assume that for any input, the size of every minimal sensitive block is at most $S$, otherwise by Lemma6.12 we get
\[D(f\comp g)\ge D(f\comp\XOR)=\Omega(S)=\Omega(\operatorname{bs}(f)^{1/5}).\]
Consider any input $x$ with maximal block sensitivity and the corresponding set of minimal sensitive blocks. The certificate complexity of $x$ is at least $9S^5$, since the certificate complexity is at least the block sensitivity. By Lemma6.14 (for $k=S$ and $K=2S^3$), we can find a substitution $\sigma$ such that $f|_\sigma$ has a parity certificate at $x|_\sigma$ of size at most $S$ with minimal weight at least $2S^3$. In addition, the block sensitivity of $f|_\sigma$ at $x|_\sigma$ is at least $9S^5-2S^4\ge7S^5$ (since we set at most $2S^4$ variables with $\sigma$), and the blocks still have size at most $S$. We fix the substitution $\sigma$ and consider the restricted function. In what follows, we denote $\hat f:=f|_\sigma$ and $\hat x:=x|_\sigma$.
Consider any input $y$ obtained from $\hat x$ by flipping one of the minimal sensitive blocks. Suppose that the certificate complexity of this input is at most $2S^3-S$. Then by Lemma6.15 (for $k=C_\oplus(\hat f,\hat x)\le S$ and $K=2S^3$), the certificate for $y$ and the parity certificate for $\hat x$ intersect. Since $\hat f(\hat x)\ne\hat f(y)$, this leads to a contradiction. So, the certificate complexity of $y$ should be greater than $2S^3-S$.
Again, by using Lemma6.14 (for $k=S$ and $K=S^2-1$), there is a substitution $\tau$ of at most $S^3$ variables, such that $y|_\tau$ has a parity certificate of size at most $S$ with minimal weight at least $S^2-1$. In addition, the certificate complexity of $y|_\tau$ is at least $2S^3-S^3=S^3$.
Since the substitution $\tau$ consists of at most $S^3$ variables, the minimal weight of the parity certificate $C_\oplus(\hat f,\hat x)|_\tau$ (we apply $\tau$ to every linear form in the certificate) is at least $2S^3-S^3\ge S^3$. Note that substitution $\tau$ can be inconsistent with $\hat x$, but we still can make this substitution inside $C_\oplus(\hat f,\hat x)$ and get a solvable certificate for $\hat f|_\tau$. Thus, if we find any point $z$, consistent with the substitution $\tau$, such that the certificate complexity of $z|_\tau$ is at most $S^2-S-1$, then this certificate $C(\hat f|_\tau,z|_\tau)$ should intersect with both $C_\oplus(\hat f,\hat x)|_\tau$ and $C_\oplus(\hat f,y)|_\tau$ by Lemma6.15 (for $k=C_\oplus(\hat x)\le S$ and $K=S^2-1$). Since all the these certificates intersect, we have $\hat f(\hat x)=\hat f(z)=\hat f(y)$, but this contradicts $\hat f(\hat x)\ne\hat f(y)$.
Now, consider an input $z$ consistent with $\tau$ such that all the $*$-variables of $\tau$ are set to zeroes. We have shown above that the certificate complexity of $\hat f|_\tau$ at $z|_\tau$ is at least $S^2-S$. Recall that $z|_\tau$ has a block sensitivity witness with all the blocks of size at most $S$. Since all the blocks together constitute a certificate, we conclude that the block sensitivity of $\hat f|_\tau$ at $z|_\tau$ is at least $S-1$. Since $z|_\tau$ is the all zeroes, by Lemma6.13 we have
\[D(f\comp g)\ge D(f\comp\AND)=\Omega(S)=\Omega(\operatorname{bs}(f)^{1/5}).\]
\end{proof}

\begin{originalresult}{Theorem 6.10 (author source p.41)}
For every gadget $g$, one of the following cases holds:
\begin{enumerate}
\item[(i)] There is an infinite family of functions $f_n$ with $\operatorname{bs}(f_n)=\Omega(n)$ and $D(f_n\comp g)=O(1)$.
\item[(ii)] For every function $f$, we have $D(f\comp g)=\Omega(\operatorname{bs}(f)^{\Omega(1)})$.
\end{enumerate}
\end{originalresult}

\begin{proof}
If $\AND,\OR\le g$, or $\AND,\XOR\le g$, or $\OR,\XOR\le g$, then by Theorem6.1 and Theorem6.9, $D(f\comp g)=\Omega(\operatorname{bs}(f)^{\Omega(1)})$, so the second case holds.
If none of these cases hold, then at least two of the functions AND, OR, XOR do not reduce to $g$. Corollary6.5 and Corollary6.6 show that either $g$ is constant (and so the first case trivially holds) or it is a blow-up of one of the functions XOR, OR, AND. By taking $f_n=\XOR_n,\OR_n,\AND_n$ (respectively), we get the first case of the theorem.
\end{proof}

\begin{thebibliography}{9}
\bibitem{zhang} Shengyu Zhang (2009). On the Tightness of the Buhrman--Cleve--Wigderson Simulation. In Proceedings of the 20th International Symposium on Algorithms and Computation,ISAAC'09,434--440. Springer. doi:10.1007/978-3-642-10631-6\_45.
\bibitem{razborov} Alexander Razborov (1992). On the distributional complexity of disjointness. Theoretical Computer Science106(2),385--390.
\bibitem{kn} Eyal Kushilevitz and Noam Nisan (1996). Communication Complexity. Cambridge University Press.
\bibitem{hatami} Hamed Hatami, Kaave Hosseini and Shachar Lovett (2018). Structure of protocols for XOR functions. SIAM Journal on Computing47(1),208--217. doi:10.1137/17M1136869.
\bibitem{zhangshi} Zhiqiang Zhang and Yaoyun Shi (2010). On the parity complexity measures of Boolean functions. Theoretical Computer Science411(26),2612--2618.
\end{thebibliography}
\end{document}
